% GENERATED FILE. Edit canonical lessons, then run npm run generate:books.
% canonical-source-hash: fnv1a64:c231f5eb1ef03194
% canonical-lessons: MR-C483-apcan, MR-C483-pitta, MR-C483-julab, MR-C483-okne, MR-C483-pacne

\chapter{An Upset Stomach}
\label{ch:mr-an-upset-stomach}
\hlchaptermodality{483}

\begin{chapteropening}
\textbf{By the end of this chapter:} \emph{I can talk about indigestion, acidity, diarrhoea and being sick.}

Complete the last lesson of chapter 483: I can talk about indigestion, acidity, diarrhoea and being sick.
\end{chapteropening}

\section[\texorpdfstring{\mr{अपचन} (apcan)}{अपचन (apcan)}]{\mr{अपचन} (apcan) — indigestion}

\label{lesson:MR-C483-apcan}



\begin{quote}
Before the new one: say the Marathi for to groan, then the Marathi for to moan in pain.
\end{quote}

\subsection*{You'll want to know: \mr{अपचन}}
\textbf{\mr{अपचन}} — \emph{apcan} — \textquotedblleft{}indigestion\textquotedblright{}.

\begin{quote}
\mr{तेलकट} \mr{खाल्ल्याने} \mr{मला} \mr{अपचन} \mr{झालं}.
\end{quote}

\begin{quote}
\emph{telkaṭ khāllyāne malā apcan jhālaṁ.}
\end{quote}

\begin{quote}
I got indigestion from eating oily food.
\end{quote}

Neuter.

\mr{अ}-, not, + \mr{पचन}, digestion. \mr{अपचन} \mr{होणे}, with the dative.

\begin{cousinweb}
\textbf{\mr{अपचन}} (\emph{apcan}) is built from Sanskrit \emph{a-}, \textquotedblleft{}not\textquotedblright{}, and \emph{pacana}, \textquotedblleft{}cooking, maturing\textquotedblright{} (Monier-Williams): a learned word for food not digested.
\end{cousinweb}

\begin{grammarlens}[title={What you've built}]
One more word for describing what you lived and felt.
\end{grammarlens}

\subsection*{Your turn}
\begin{itemize}
\raggedright
  \item \emph{Say it:} \emph{apcan}
  \item \emph{Say it:} \emph{apcan}, once more
  \item \emph{Say it:} \emph{telkaṭ khāllyāne malā apcan jhālaṁ.} — I got indigestion from eating oily food.
  \item \emph{Recall:} say the Marathi for to groan, then the Marathi for to moan in pain, then say \emph{apcan} again
\end{itemize}

\begin{tcolorbox}[breakable,colback=teal!4,colframe=teal!35!black,title={Before you move on}]
What does \mr{अपचन} mean? (\textquotedblleft{}Indigestion\textquotedblright{}.) Say it once more.
\end{tcolorbox}
\section[\texorpdfstring{\mr{पित्त} (pitta)}{पित्त (pitta)}]{\mr{पित्त} (pitta) — acidity; bile}

\label{lesson:MR-C483-pitta}



\begin{quote}
Before the new one: say the Marathi for to moan in pain, then the Marathi for indigestion.
\end{quote}

\subsection*{You'll want to know: \mr{पित्त}}
\textbf{\mr{पित्त}} — \emph{pitta} — \textquotedblleft{}acidity; bile\textquotedblright{}.

\begin{quote}
\mr{उपाशी} \mr{राहिल्यामुळे} \mr{त्याला} \mr{पित्त} \mr{झालं}.
\end{quote}

\begin{quote}
\emph{upāśī rāhilyāmuḷe tyālā pitta jhālaṁ.}
\end{quote}

\begin{quote}
He got acidity from going without food.
\end{quote}

Neuter.

\mr{पित्त} \mr{होणे}, to get acidity, with the dative.

\begin{cousinweb}
\textbf{\mr{पित्त}} (\emph{pitta}) is a tatsama, taken from Sanskrit \emph{pitta}, \textquotedblleft{}bile\textquotedblright{} (CDIAL 8181).
\end{cousinweb}

\begin{grammarlens}[title={What you've built}]
One more word for describing what you lived and felt.
\end{grammarlens}

\subsection*{Your turn}
\begin{itemize}
\raggedright
  \item \emph{Say it:} \emph{pitta}
  \item \emph{Say it:} \emph{pitta}, once more
  \item \emph{Say it:} \emph{upāśī rāhilyāmuḷe tyālā pitta jhālaṁ.} — He got acidity from going without food.
  \item \emph{Recall:} say the Marathi for to moan in pain, then the Marathi for indigestion, then say \emph{pitta} again
\end{itemize}

\begin{tcolorbox}[breakable,colback=teal!4,colframe=teal!35!black,title={Before you move on}]
What does \mr{पित्त} mean? (\textquotedblleft{}Acidity; bile\textquotedblright{}.) Say it once more.
\end{tcolorbox}
\section[\texorpdfstring{\mr{जुलाब} (julāb)}{जुलाब (julāb)}]{\mr{जुलाब} (julāb) — diarrhoea}

\label{lesson:MR-C483-julab}



\begin{quote}
Before the new one: say the Marathi for indigestion, then the Marathi for acidity.
\end{quote}

\subsection*{You'll want to know: \mr{जुलाब}}
\textbf{\mr{जुलाब}} — \emph{julāb} — \textquotedblleft{}diarrhoea\textquotedblright{}.

\begin{quote}
\mr{शिळं} \mr{अन्न} \mr{खाल्ल्याने} \mr{बाळाला} \mr{जुलाब} \mr{झाले}.
\end{quote}

\begin{quote}
\emph{śiḷaṁ anna khāllyāne bāḷālā julāb jhāle.}
\end{quote}

\begin{quote}
The baby got diarrhoea from eating stale food.
\end{quote}

Masculine.

Usually plural: \mr{जुलाब} \mr{होणे}, with the dative.

\begin{cousinweb}
\textbf{\mr{जुलाब}} (\emph{julāb}) is a loan from Persian \emph{julāb}, \textquotedblleft{}a purgative drink\textquotedblright{}, usually traced to Arabic.
\end{cousinweb}

\begin{grammarlens}[title={What you've built}]
One more word for describing what you lived and felt.
\end{grammarlens}

\subsection*{Your turn}
\begin{itemize}
\raggedright
  \item \emph{Say it:} \emph{julāb}
  \item \emph{Say it:} \emph{julāb}, once more
  \item \emph{Say it:} \emph{śiḷaṁ anna khāllyāne bāḷālā julāb jhāle.} — The baby got diarrhoea from eating stale food.
  \item \emph{Recall:} say the Marathi for indigestion, then the Marathi for acidity, then say \emph{julāb} again
\end{itemize}

\begin{tcolorbox}[breakable,colback=teal!4,colframe=teal!35!black,title={Before you move on}]
What does \mr{जुलाब} mean? (\textquotedblleft{}Diarrhoea\textquotedblright{}.) Say it once more.
\end{tcolorbox}
\section[\texorpdfstring{\mr{ओकणे} (okṇe)}{ओकणे (okṇe)}]{\mr{ओकणे} (okṇe) — to vomit; to throw up}

\label{lesson:MR-C483-okne}



\begin{quote}
Before the new one: say the Marathi for acidity, then the Marathi for diarrhoea.
\end{quote}

\subsection*{You'll want to know: \mr{ओकणे}}
\textbf{\mr{ओकणे}} — \emph{okṇe} — \textquotedblleft{}to vomit; to throw up\textquotedblright{}.

\begin{quote}
\mr{बसमध्ये} \mr{बाळाने} \mr{दूध} \mr{ओकलं}.
\end{quote}

\begin{quote}
\emph{basmadhye bāḷāne dūdh oklaṁ.}
\end{quote}

\begin{quote}
The baby threw up its milk on the bus.
\end{quote}

Past, with -\mr{ने} (he …): \mr{त्याने} \mr{ओकलं}, \emph{tyāne oklaṁ}.

Blunter than \mr{उलटी} \mr{होणे}; what comes up is the object.

\begin{cousinweb}
\textbf{\mr{ओकणे}} (\emph{okṇe}) is inherited: CDIAL lists Marathi \emph{okṇẽ}, \textquotedblleft{}vomit\textquotedblright{}, under the expressive \emph{*ōkk} (CDIAL 2538).
\end{cousinweb}

\begin{grammarlens}[title={What you've built}]
One more word for describing what you lived and felt.
\end{grammarlens}

\subsection*{Your turn}
\begin{itemize}
\raggedright
  \item \emph{Say it:} \emph{okṇe}
  \item \emph{Say it:} \emph{okṇe}, once more
  \item \emph{Say it:} \emph{basmadhye bāḷāne dūdh oklaṁ.} — The baby threw up its milk on the bus.
  \item \emph{Recall:} say the Marathi for acidity, then the Marathi for diarrhoea, then say \emph{okṇe} again
\end{itemize}

\begin{tcolorbox}[breakable,colback=teal!4,colframe=teal!35!black,title={Before you move on}]
What does \mr{ओकणे} mean? (\textquotedblleft{}To vomit; to throw up\textquotedblright{}.) Say it once more.
\end{tcolorbox}
\section[\texorpdfstring{\mr{पचणे} (pacṇe)}{पचणे (pacṇe)}]{\mr{पचणे} (pacṇe) — to be digested}

\label{lesson:MR-C483-pacne}



\begin{quote}
Before the new one: say the Marathi for diarrhoea, then the Marathi for to vomit.
\end{quote}

\subsection*{You'll want to know: \mr{पचणे}}
\textbf{\mr{पचणे}} — \emph{pacṇe} — \textquotedblleft{}to be digested\textquotedblright{}.

\begin{quote}
\mr{रात्रीचं} \mr{जड} \mr{जेवण} \mr{लवकर} \mr{पचत} \mr{नाही}.
\end{quote}

\begin{quote}
\emph{rātrīcaṁ jaḍ jevaṇ lavkar pacat nāhī.}
\end{quote}

\begin{quote}
A heavy dinner does not digest quickly.
\end{quote}

Past (to him …): \mr{त्याला} \mr{पचलं}, \emph{tyālā paclaṁ}.

The eater takes the dative: \mr{मला} \mr{पचलं}. \mr{पचवणे}, to digest.

\begin{cousinweb}
\textbf{\mr{पचणे}} (\emph{pacṇe}) is a tadbhava, inherited from Sanskrit \emph{pacyate}; CDIAL lists Marathi \emph{pacṇẽ}, \textquotedblleft{}to be digested\textquotedblright{} (CDIAL 7654).
\end{cousinweb}

\begin{grammarlens}[title={What you've built}]
One more word for describing what you lived and felt.
\end{grammarlens}

\subsection*{Your turn}
\begin{itemize}
\raggedright
  \item \emph{Say it:} \emph{pacṇe}
  \item \emph{Say it:} \emph{pacṇe}, once more
  \item \emph{Say it:} \emph{rātrīcaṁ jaḍ jevaṇ lavkar pacat nāhī.} — A heavy dinner does not digest quickly.
  \item \emph{Recall:} say the Marathi for diarrhoea, then the Marathi for to vomit, then say \emph{pacṇe} again
\end{itemize}

\begin{tcolorbox}[breakable,colback=teal!4,colframe=teal!35!black,title={Before you move on}]
What does \mr{पचणे} mean? (\textquotedblleft{}To be digested\textquotedblright{}.) Say it once more.
\end{tcolorbox}
